\subsection{Repairing the actual formal pair generators}
\label{subsec:sr-source-frame-truncation}

All rings in this subsection are arbitrary commutative
\(\mathbb Q\)-algebras.  In particular, none of the
coefficient comparisons below assumes a reduced ring.
Fix the source packet \(E,H,D_w,l,X_o,c\), the equality
\(\star a=H.\mathrm{bar}(a)\) on \(\mathcal O_E\), a
\(c.F\)-algebra \(S\), an actual framed point \(y\), and
an actual \(C:\mathrm{ChartAt}(X_o,c,y)\).
Its formal group is \(G\), its coefficient actions are
\(A_a\), and its pair generators are \(e_m\).
Write
\[
 J_m^{(s)}=S[Z_{b,i}:1\le b\le s,\ 1\le i\le d]/
                         (Z)^{m+1},\qquad j_m f=f\bmod(Z)^{m+1}.
 \tag{FC1}
\]
Here \(d=C.d\), and every point, line generator and
comparison in the argument belongs to this same chart.

\begin{srlemma}[The normalized logarithm used by the correction]
\label{lem:sr-source-normalized-formal-log}
There is a formal logarithm \(\ell_G(X)=X+O((X)^2)\) with
\[
 \ell_G(G(X,Y))=\ell_G(X)+\ell_G(Y),\qquad
 \ell_G(A_aX)=A_a^{(1)}\ell_G(X),\quad A_a^{(1)}=dA_a(0).
 \tag{FC2}
\]
A homomorphism from \(G\) to the additive formal group
with zero constant term is determined by its linear term.
These assertions hold coefficientwise over \(S\).
\end{srlemma}
\begin{proof}
We give the source's invariant-differential construction.
Let \(J(X)=\partial_YG(X,0)\), so \(J(0)=I\), and put
\(\Omega(X)=J(X)^{-1}\).
Associativity, differentiated in its final variable at
zero, gives
\[
 J(G(X,Y))=\partial_YG(X,Y)J(Y).
\]
Inverting \(J\) gives
\(\Omega(G(X,Y))\partial_YG(X,Y)=\Omega(Y)\).
Commutativity gives the corresponding identity with
\(X\) and \(Y\) exchanged.
Differentiating these identities once more, commuting
the two partial derivatives and multiplying by the
invertible \(J(X)\), gives
\[
 \sum_jJ_{jb}\partial_a\Omega_{ij}
       =\sum_j(\partial_j\Omega_{ia})J_{jb}.
\]
Multiply on the right by \(J^{-1}\) to obtain
\(\partial_a\Omega_{ij}=\partial_j\Omega_{ia}\).
Thus every row of \(\Omega\) is a closed formal one-form.

For a closed form \(\omega=(\omega_i)\), define its
zero-constant primitive by the following actual coefficient:
\[
 [X^\mu]\mathrm{prim}(\omega)=
 \begin{cases}
 |\mu|^{-1}\displaystyle\sum_{i:\mu_i>0}
                    [X^{\mu-e_i}]\omega_i,&|\mu|>0,\\
 0,&\mu=0.
 \end{cases}                                                   \tag{FC3}
\]
Closedness and \(\sum_i\mu_i=|\mu|\) show, coefficient
by coefficient, that
\(\partial_i\mathrm{prim}(\omega)=\omega_i\).
For example the coefficient at \(X^\nu\) in the left
side is \((\nu_i+1)[X^{\nu+e_i}]\mathrm{prim}(\omega)\);
closedness changes the summand indexed by \(r\) into
\((\nu_r+\delta_{ri})[X^\nu]\omega_i/(\nu_i+1)\),
and the sum gives the claimed coefficient.
All positive integers used here are units in \(S\).
Set \(\ell_{G,i}=\mathrm{prim}(\Omega_{i,*})\).
Then \(d\ell_G=\Omega\) and \(d\ell_G(0)=I\).
The preceding matrix identities make every partial
derivative of
\(\ell_G(G(X,Y))-\ell_G(X)-\ell_G(Y)\) zero.
A formal series over a \(\mathbb Q\)-algebra with all
partial derivatives zero is constant; evaluate at the
origin to obtain the first identity in (FC2).

More generally, if \(h(G(X,Y))=h(X)+h(Y)\), differentiation
in \(Y\) at zero gives
\(dh(X)J(X)=dh(0)\).
Thus its differential and, with \(h(0)=0\), all its
coefficients are determined by \(dh(0)\).
The maps \(\ell_GA_a\) and \(A_a^{(1)}\ell_G\) are
homomorphisms with the same differential, proving the
second identity in (FC2).
This is also the source's formal-homomorphism uniqueness
argument.
\end{proof}

For a matrix \(B\) over \(S\), put
\[
 {\cal B}_B(X,Y)=\ell_G(X)^tB\ell_G(Y),\qquad
 T_B(X,Y)=\exp(-{\cal B}_B(X,Y)).
 \tag{FC4}
\]
The exponential is a formal series: its argument has
zero constant term, and its reduction to any \(J_m^{(2)}\)
is a finite sum.  Its inverse is \(\exp({\cal B}_B)\).

\begin{srlemma}[Classification of a finite pair discrepancy]
\label{lem:sr-source-pair-discrepancy-exponential}
Let \(U\in(J_m^{(2)})^\times\) satisfy, as identities in
the indicated truncated polynomial rings,
\[
 \begin{split}
 U(G(X,Y),Z)&=U(X,Z)U(Y,Z),\\
 U(Y,X)&=U(X,Y),\\
 U(A_aX,Y)&=U(X,A_{\bar a}Y)\qquad(a\in\mathcal O_E).
 \end{split}                                                   \tag{FC5}
\]
There is a matrix \(b\) with
\[
 b^t=b,\qquad (A_a^{(1)})^tb=bA_{\bar a}^{(1)},\qquad
                         U=j_mT_b.                           \tag{FC6}
\]
For \(m=0,1\) one may take \(b=0\).
\end{srlemma}
\begin{proof}
Choose a polynomial representative \(P\) of \(U\).
Its constant coefficient \(u\) is a unit, by applying
the augmentation to \(U\).
The first identity in (FC5) gives \(u=u^2\), hence \(u=1\).

Here is the leading-term argument used by the source.
If \(P-1\) has no terms of degree below \(j\), its
degree-\(j\) component \(H_j(X,Z)\), for \(1\le j\le m\),
satisfies
\[
 H_j(X+Y,Z)=H_j(X,Z)+H_j(Y,Z),\qquad H_j(Z,X)=H_j(X,Z).
 \tag{FC7}
\]
Indeed \(G(X,Y)=X+Y+O((X,Y)^2)\), and all products
of two positive residual terms have degree at least
\(2j>j\).
The second assertion makes \(H_j\) additive in its
second variable as well.
An additive polynomial over a \(\mathbb Q\)-algebra
is homogeneous of degree one in that variable:
put \(Y=X\) in its additivity identity and compare
homogeneous degrees; \(2^r-2\) is a nonzero rational
unit for \(r\ne1\).
Thus a polynomial additive separately in \(X,Z\) has
bidegree \((1,1)\).
For \(j=1\) this forces \(H_j=0\).
For \(j=2\) it gives \(H_2=X^tBZ\), with \(B^t=B\).
If \(m<2\), the same argument already proves \(U=1\).

Suppose \(m\ge2\).
Comparison of the degree-two terms in the last identity
of (FC5) gives
\[
                 (A_a^{(1)})^tB=BA_{\bar a}^{(1)}.            \tag{FC8}
\]
By (FC2), \({\cal B}_B\) is additive with respect to \(G\)
in either argument; it is symmetric and has the same
adjoint balance.
Consequently \(j_m\exp(-{\cal B}_B)\) satisfies (FC5).
It begins \(1-X^tBY+O((X,Y)^3)\).
The product
\(V=Uj_m\exp(-{\cal B}_B)\) therefore still satisfies
(FC5) and has no residual term of degree below three.
For \(j=3,\ldots,m\), apply (FC7) successively:
a separately additive degree-\(j\) component must be
zero because it would have total degree two.
Thus \(V=1\).
It follows that \(U=j_m\exp({\cal B}_B)=j_mT_{-B}\).
Take \(b=-B\).
This proof compares coefficients and cancels units only;
nilpotent coefficients cause no additional case.
\end{proof}

\begin{srproposition}[Actual compatible pair generators]
\label{prop:sr-source-actual-frame-truncation-fix}
For the actual chart \(C\), there are generators \(e'_m\)
of its same pair lines and matrices \(C_m\) such that
\[
 \begin{gathered}
 C_0=0,\quad C_m^t=C_m,\quad
 (A_a^{(1)})^tC_m=C_mA_{\bar a}^{(1)},\\
                 e'_m=(j_mT_{C_m})e_m.                       \tag{FC9}
 \end{gathered}
\]
The corrected generators satisfy the original first
and second partial laws, symmetry, adjointness and
origin rigidification, with the same structure series
as \(C\); in addition they agree under every truncation.
This constructs the literal \(\mathrm{FrTruncFixSt}\).
\end{srproposition}
\begin{proof}
First construct the actual discrepancy; it is not a
freely chosen scalar.
The chart point identities identify the pullback of
\(e_{m+1}\) to level \(m\) with the same pair line as
\(e_m\).  Two generators of the same line differ by
one global unit \(U_m\in(J_m^{(2)})^\times\):
on a line chart take their ratio, and use the common
transition function to glue these local unit ratios.
Thus
\[
                       \rho_m e_{m+1}=U_me_m.               \tag{FC10}
\]
This construction is on the entire scheme
\(\operatorname{Spec}J_m^{(2)}\).

Compare the first partial law at levels \(m+1\) and \(m\).
The triple chart points agree after truncation, by the
actual point law.  Pullback of the fixed cubical line
generator also agrees.
The two truncations of the same structure series
\(\mathrm{ph1}\) agree coefficientwise.
In each stalk cancel these common unit factors and
the three pair generators.  What remains is exactly
\[
                U_m(G(X,Y),Z)=U_m(X,Z)U_m(Y,Z).
\]
The same comparison of symmetry and adjointness gives
the other two identities (FC5).
Equality in every stalk gives equality of the global
units.  Comparing complex points alone would not
justify this step for the allowed \(S\).

Apply Lemma~\ref{lem:sr-source-pair-discrepancy-exponential}
to obtain symmetric balanced \(b_m\) with \(U_m=j_mT_{b_m}\).
Define recursively, with the source's signs,
\[
                         C_0=0,\qquad C_{m+1}=C_m-b_m.       \tag{FC11}
\]
Symmetry and balance survive subtraction.
Set \(e'_m\) by (FC9).  Then
\[
 \rho_m e'_{m+1}
 =j_mT_{C_{m+1}}U_me_m
 =j_mT_{C_{m+1}+b_m}e_m=e'_m.                               \tag{FC12}
\]
For example the multiplier in the first partial law is
\[
 T_{C_m}(G(X,Y),Z)\,
              T_{C_m}(X,Z)^{-1}T_{C_m}(Y,Z)^{-1}=1.
\]
The multiplier for the second partial law is \(1\)
by additivity in the second argument.
The symmetry multiplier is \(1\) since \(C_m\) is
symmetric; the adjoint multiplier is \(1\) by its
balance and (FC2).
Thus the same four original structure series still
witness the four frame laws, not new unrelated series.
At level zero \(C_0=0\), so the original rigidification
is unchanged.  This proves all the stated clauses.
\end{proof}

The source next needs actual section germs for these
corrected generators.  The following inverse-limit
argument supplies them with their precise finite-order
change, even though the matrices \(C_m\) themselves
need not form a compatible inverse system.
Define
\[
 Q_B(Z)=\ell_G(Z)^tB\ell_G(A_{w_0}Z).                        \tag{FC13}
\]
This is the source \(\mathrm{diagL}\).
The single-line generator is the primitive generator
multiplied by the pair generator on the graph
\((Z,A_{w_0}Z)\); hence (FC9) multiplies it by
\(j_m\exp(-Q_{C_m})\).

\begin{srproposition}[The actual corrected germ and its sign]
\label{prop:sr-source-actual-frame-germ-fix}
Let \(k\ge0\), let \(F\) be an actual section of the
specified \(k\)-th power of the polarizing line, and
let \(f\) be its actual germ in \(C\).
There is an actual germ \(f'\) of the same \(F\)
for the corrected chart, and
\[
              j_Df'=j_D\bigl(f\exp(kQ_{C_D})\bigr)
                                      \qquad(D\ge0).        \tag{FC14}
\]
Thus \(\mathrm{GermFixSt}\) is constructed, with its
quantifier over every section and every order \(D\).
\end{srproposition}
\begin{proof}
We spell out the source's general germ construction for
a truncation-compatible frame.
Pull \(F\) to each actual jet base.
The corrected primitive-times-pair generator is a
generator of the actual single line there: its tensor
factors cancel by the actual chart identity, including
the \(A_{w_0}\)-graph.
Raise it to the \(k\)-th power.
On any common line chart divide the pulled section
coefficient by this unit generator.
The transition factors cancel, so these local functions
glue to one global value \(v_m\in\Gamma(\operatorname{Spec}
J_m^{(1)},\mathcal O)=J_m^{(1)}\).
The section pullbacks, chart points and corrected
generators all agree under truncation.
In every stalk their value equations are therefore the
same equation with the same unit generator.
Cancel the unit and use equality of sections detected
in stalks.  This gives \(\rho_mv_{m+1}=v_m\).

Choose polynomial representatives \(p_m\) of \(v_m\).
Their coefficients of degree at most \(m\) are independent
of later representatives, since
\(p_l-p_m\in(Z)^{m+1}\) for \(l\ge m\).
Define
\([Z^\mu]f'=[Z^\mu]p_{|\mu|}\).
Then \(j_mf'=v_m\) for every \(m\).
The section pullbacks and generators used to define
\(v_m\) give exactly the witnesses in the source
definition of \(\mathrm{IsGerm}\).

To compare with \(f\), use its supplied actual
pullback-section and value witnesses at level \(D\).
Equation (FC9) changes its single-line generator by
\(u_D=j_D\exp(-Q_{C_D})\), and its \(k\)-th power
generator by \(u_D^k\).
The coordinate of the same section consequently changes
by \(u_D^{-k}\).
Since
\[
 j_D\exp(kQ_{C_D})\,u_D^k=1,
\]
this coordinate is the right side of (FC14).
Uniqueness of section pullback, equality of generators
expressed by \(\mathrm{Corr}\), and uniqueness of a
value in a unit generator identify it with \(v_D\).
The exponent in (FC14) is positive; it is forced by
division by the corrected generator.
\end{proof}

\begin{srlemma}[The actual eigenvectors and the balanced support]
\label{lem:sr-source-balanced-correction-conjugate-support}
For a source letter \(\lambda\in\mathrm{Lt1}(D_w)\),
write \(\tau_\lambda\) for its underlying embedding.
Conjugation in the support assertion compares these embeddings,
with no assertion about the multiplicity index.
For every actual chart as above,
\[
 A_a^{(1)}u_\lambda
    =\mathrm{embK}_c(\tau_\lambda(a))u_\lambda.                  \tag{FC15}
\]
Consequently a balanced matrix \(B\) satisfies
\[
 u_\lambda^tBu_{\lambda'}\ne0
                  \quad\Longrightarrow\quad
                          \tau_{\lambda'}=\overline{\tau_\lambda}.           \tag{FC16}
\]
These are identities over \(S\), including its zero divisors.
\end{srlemma}
\begin{proof}
Let \(D_i\) be the actual dual-number tangent point
obtained from the \(i\)-th coordinate of the chart.
The chart's tangent law says that the combination
\(\sum_i u_{\lambda,i}D_i\) is the supplied framed
eigentangent \(y.u_\lambda\).
Differentiating the actual chart action law makes the
action on these coordinate combinations the matrix
\(A_a^{(1)}\).

Here is the source's injectivity check for this tangent
comparison.  The chart basis clause writes any coordinate
vector uniquely as a combination of the \(u_\lambda\).
Group its tangent combination by its embedding \(\lambda\).
Distinct embeddings of \(E\) differ on some
\(a\in\mathcal O_E\): write any field element as a ratio
of integral elements to see that equality on
\(\mathcal O_E\) would imply equality on \(E\).
Their nonzero difference in \(c.F\) maps to a unit of \(S\).
If a finite sum of eigentangents is zero, apply
\([a]-\lambda_2(a)\) to remove one component and induct
on the finite set of embeddings; the unit differences
separate the remaining component.
It follows that each grouped eigentangent is zero.
The actual basis clause of \(y.u\) within that eigenspace
then makes all its coefficients zero.
Together with the chart basis clause this makes the
coordinate tangent map injective.
The two expressions for the action on \(y.u_\lambda\)
therefore give the coordinate equality (FC15).

For \(b=u_\lambda^tBu_{\lambda'}\), balance and (FC15)
give
\[
 \bigl(\mathrm{embK}_c(\tau_\lambda(a))-
       \mathrm{embK}_c(\tau_{\lambda'}(\bar a))\bigr)b=0
                                      \quad(a\in\mathcal O_E).
\]
If \(\tau_{\lambda'}\ne\overline{\tau_\lambda}\), integral embedding
separation gives an \(a\) for which the factor is a
nonzero element of \(c.F\), hence a unit after passage
to \(S\).
Cancel it to obtain \(b=0\).
This proves (FC16) by unit cancellation, without a
domain hypothesis on \(S\).
\end{proof}

\paragraph{Frozen source correspondence and exact remaining scope.}
The normalized-logarithm construction has prefix
\texttt{src/ReciprocityTheory/ReductiveGroups/ArithmeticGeometry/}
\texttt{TaylorDiscriminants/WickVanishing/}.
It follows \texttt{LogInv.lean}
(\texttt{0c704c3b629b0bde3f30457a87951cacbe00fe15}),
\texttt{LogClosed.lean}
(\texttt{364e6ec267ab62e6ebc16173142d239e0025a497}),
\texttt{LogPoinc.lean}
(\texttt{9d5c0658df322bbf5ff0ca25e38651b6a6ec953b}),
and \texttt{PderivSecondEvFirstPointEqZero.lean}
(\texttt{b7fd0380bc4ed3c9283782a057debc08adc0818e}).
Formula (FC3) is literally \(\mathrm{prim}\);
(FC2) is \(\mathrm{logOf\_add}\).
The homomorphism uniqueness is the coefficient argument
of \texttt{HomUniq.lean}
(\texttt{0dd9e559e8470dee9c55401a59524e6929947919}).

The remaining files have prefix
\texttt{src/Formalism/SharpRestriction/Divisibility/GeneratorPowers/}.
The actual discrepancy (FC10), including its stalkwise
three comparisons, is \texttt{FrRatio.lean}
(\texttt{00628f81e6a1248458637030d61594dd5f1a70f1}),
\(\mathrm{exists\_ratioU}\), \(\mathrm{frRatio\_holds}\).
The leading homogeneous argument and balance are
\texttt{FormalGroupLinearSecond.lean}
(\texttt{75033f98de459ded476de15ead857773d6612052}),
\texttt{BalOf.lean}
(\texttt{ae3e2a62fc1bcf4b76a336e9506551a26a4a1f6e}),
and \texttt{IsUnitCoeffZero.lean}
(\texttt{f8fc4377b7b709bd4ddc0207c7a8697f5f77430e}),
\(\mathrm{fixExp}\).
The source's actual logarithm is \(X+O(X^2)\), as
verified in \texttt{LogLowOrderOne.lean}
(\texttt{721a26bb1ff0b12ba6596e60cfaef6659454b43e}).

Preservation of the two partial laws and of symmetry
and adjointness is
\texttt{ConstantCoeffAddSecondThird.lean}
(\texttt{91009c38cc83a1e0bc139b88ccbb02a61b1b8cfd})
and \texttt{ConstantCoeffAddFirstSecond.lean}
(\texttt{b1d210e3eb794f0d359aca47f885bcb20dfb0a76}),
\(\mathrm{fixTw}\).
The recursion (FC11) is \(\mathrm{ccOf}\) in
\texttt{CIndFrFixAsm.lean}
(\texttt{4155678eb405adf4d0539fc6f2322765e7644a20}).

The actual pullback/value construction and compatible
coefficients are \texttt{GermDVal.lean}
(\texttt{a45f227868b66b0d323a9ac91ab71631d20059a6}),
\texttt{GermDJet.lean}
(\texttt{0515703d03f57db59ec97f20d43c74eb03193865}),
and \texttt{GermDPf.lean}
(\texttt{fc6c753be8c9cdcb999ecea9bb07ccd4a8096977}),
\(\mathrm{germSt\_holds}\).
The exact coordinate correction (FC14) is
\(\mathrm{germ\_order}\), \(\mathrm{germFix\_holds}\)
in \texttt{EvalSinglePointId.lean}
(\texttt{228b49fb832031923cb127519a268231f5439ad0}).
The actual tangent separation is \texttt{LinEig.lean}
(\texttt{523682e6ce5c5fa038ad07cce17a8720650d2436});
(FC16) is \(\mathrm{pair\_of\_balanced}\) in
\texttt{CIndVerbB.lean}
(\texttt{648a74827aba4ad62881a418f592ec4e62f56390}).

This constructs the actual truncation correction and
the corresponding actual section germs.
Full comparison of two different corrected charts still
requires their exponential-coordinate identification,
the global primitive-line comparison, and harmonic
annihilation of the resulting conjugate quadratic term.
Those are the remaining branches of
\(\mathrm{chartIndep2\_of\_glob}\); this subsection does
not replace them by a formal assertion of invariance.
